\chapter{Eigenvalues and Eigenvectors}
\label{ch:i09-eigenvalues-and-eigenvectors}

Eigenvectors are directions preserved by an operator up to scaling, and eigenvalues are the corresponding scale factors. The characteristic polynomial packages all eigenvalues algebraically. This chapter emphasizes both geometric meaning and the limitations of eigenvector methods when a full eigenbasis does not exist.

\section*{Learning goals}
After this chapter, the reader should be able to:
\begin{itemize}
\item compute eigenvalues from the characteristic polynomial and compute a basis for each eigenspace;
\item distinguish algebraic multiplicity from geometric multiplicity in explicit examples;
\item prove that eigenvectors belonging to distinct eigenvalues are linearly independent;
\item decide whether a given vector or subspace is preserved by an operator using eigenvector information;
\item construct examples with repeated eigenvalues that are or are not diagonalizable;
\item check an eigenpair directly and diagnose arithmetic errors by substituting back into \(Tv=\lambda v\);
\end{itemize}
\section*{Conceptual roadmap}
The chapter is organized around a repeated pattern:
\[
\boxed{\text{definition}\;\longrightarrow\;\text{structural theorem}\;\longrightarrow\;\text{coordinate calculation}\;\longrightarrow\;\text{application}.}
\]
Whenever possible, we separate the invariant mathematical object from a chosen coordinate representation.

\section{Eigenvectors and eigenvalues}
A nonzero $v$ is an eigenvector of $T$ with eigenvalue $\lambda$ when $T(v)=\lambda v$.


% BEGIN VOL01-EXPANSION I09-example-01
\begin{example}[A complete eigensystem calculation]\label{ex:i09-expand-01}
For
\[
A=\begin{pmatrix}4&1\\0&2\end{pmatrix},
\]
the characteristic equation is
\((4-\lambda)(2-\lambda)=0\), so \(\lambda=4,2\).
For \(\lambda=4\), \((A-4I)v=0\) gives \(y=0\), hence
\(E_4=\operatorname{span}(1,0)\). For \(\lambda=2\),
\(2x+y=0\), so \(E_2=\operatorname{span}(1,-2)\).

Check:
\[
A(1,0)^T=4(1,0)^T,\qquad
A(1,-2)^T=2(1,-2)^T.
\]
The substitution check catches sign errors that the characteristic polynomial
alone cannot reveal.
\end{example}
% END VOL01-EXPANSION I09-example-01
\section{Eigenspaces}
$E_\lambda=\ker(T-\lambda I)$ is a subspace. Distinct eigenspaces meet only at zero.

\section{Characteristic polynomial}
For a matrix $A$, $\chi_A(t)=\det(tI-A)$. Its roots are precisely the eigenvalues over the chosen field.


% BEGIN VOL01-EXPANSION I09-example-02
\begin{example}[Field dependence of the spectrum]\label{ex:i09-expand-02}
For the quarter-turn matrix
\[
R=\begin{pmatrix}0&-1\\1&0\end{pmatrix},
\]
\[
\chi_R(t)=t^2+1.
\]
There are no real eigenvalues, but over \(\mathbb C\) the roots are
\(\pm i\). Solving \((R-iI)v=0\) gives, for example,
\(v=(1,-i)\).

The operator has not changed; extending the scalar field has supplied roots
that were unavailable over \(\mathbb R\).
\end{example}
% END VOL01-EXPANSION I09-example-02
\section{Algebraic and geometric multiplicity}
The algebraic multiplicity is the multiplicity of a root of $\chi_A$; the geometric multiplicity is $\dim E_\lambda$ and never exceeds it.


% BEGIN VOL01-EXPANSION I09-example-03
\begin{example}[Repeated eigenvalue: diagonalizable or defective?]\label{ex:i09-expand-03}
Compare
\[
I=\begin{pmatrix}1&0\\0&1\end{pmatrix},
\qquad
J=\begin{pmatrix}1&1\\0&1\end{pmatrix}.
\]
Both have characteristic polynomial \((t-1)^2\). For \(I\),
\(E_1=\mathbb R^2\), so geometric multiplicity is \(2\). For \(J\),
\((J-I)v=0\) forces \(y=0\), so geometric multiplicity is \(1\).

Thus eigenvalues with algebraic multiplicity do not determine similarity:
eigenspace dimension records additional structure.
\end{example}
% END VOL01-EXPANSION I09-example-03
\section{Independence of eigenvectors}
Eigenvectors associated with distinct eigenvalues are linearly independent.

\section{Field dependence}
A real matrix may have no real eigenvalues but acquire complex eigenvalues after extending scalars to $\mathbb C$.

\section{Core structural results}

\begin{theorem}[Distinct eigenvalues give independence]\label{thm:i09-01}
Eigenvectors corresponding to distinct eigenvalues are linearly independent.
\begin{proof}
Use induction. Apply $T-\lambda_r I$ to a dependence relation; the last term vanishes and the remaining coefficients are multiplied by nonzero scalars.
\end{proof}
\end{theorem}

\begin{theorem}[Geometric multiplicity bound]\label{thm:i09-02}
For an eigenvalue $\lambda$, $1\le\dim E_\lambda\le$ its algebraic multiplicity.
\begin{proof}
The lower bound is the definition of eigenvalue. The upper bound follows after choosing a basis beginning with an eigenspace basis, which forces a block with $\lambda$ repeated on the diagonal.
\end{proof}
\end{theorem}

\section{Worked examples}

\begin{example}[Diagonal matrix]\label{ex:i09-worked-01}
For $D=\operatorname{diag}(d_1,\ldots,d_n)$, each standard vector is an eigenvector with eigenvalue $d_i$.
\end{example}

\begin{example}[Rotation]\label{ex:i09-worked-02}
A nontrivial planar rotation has no real eigenvectors unless the angle is $0$ or $\pi$, but over $\mathbb C$ it has eigenvalues $e^{\pm i\theta}$.
\end{example}

\begin{example}[Shear]\label{ex:i09-worked-03}
$\begin{pmatrix}1&1\\0&1\end{pmatrix}$ has only the eigenvalue $1$ and only a one-dimensional eigenspace.
\end{example}

\section{Diagnostic checklist}
Before moving on, it is useful to distinguish four levels of understanding:
\begin{enumerate}
\item recognize the definitions in unfamiliar examples;
\item prove the core structural statements without coordinates when possible;
\item translate the same statement into a matrix or coordinate computation;
\item know which hypotheses are essential and produce a counterexample when one is removed.
\end{enumerate}

% BEGIN VOLUME01 SOLVED DOSSIERS
\section{Solved dossiers}
These dossiers are longer worked problems. They complement the shorter exercise--hint--solution triads by combining several ideas in one argument.

\begin{problem}[Defective shear]\label{prob:i09-dossier-01}
For $A=\begin{pmatrix}1&1\\0&1\end{pmatrix}$ compute algebraic and geometric multiplicities of $\lambda=1$.
\end{problem}
\begin{solution}
The characteristic polynomial is $(t-1)^2$, so algebraic multiplicity is $2$. The eigenspace solves $(A-I)v=0$, hence $y=0$, so it is one-dimensional. Geometric multiplicity $1<2$, making the matrix defective.
\end{solution}

\begin{problem}[Rotation and complexification]\label{prob:i09-dossier-02}
Find the eigenvalues of the real rotation matrix $\begin{pmatrix}0&-1\\1&0\end{pmatrix}$ over $\mathbb R$ and $\mathbb C$.
\end{problem}
\begin{solution}
The characteristic polynomial is $t^2+1$. It has no real roots, so no real eigenvalues; over $\mathbb C$ the eigenvalues are $i$ and $-i$.
\end{solution}

\begin{problem}[Parameter collision]\label{prob:i09-dossier-03}
For $A_t=\begin{pmatrix}t&1\\0&2\end{pmatrix}$, when are there two distinct eigenvalues?
\end{problem}
\begin{solution}
The eigenvalues are the diagonal entries $t$ and $2$. They are distinct exactly when $t\ne2$. At $t=2$ the matrix becomes a nontrivial Jordan-type block with only one eigendirection.
\end{solution}

\begin{problem}[Eigenvectors of powers]\label{prob:i09-dossier-04}
If $Tv=\lambda v$, what is $T^kv$?
\end{problem}
\begin{solution}
Induction gives $T^kv=\lambda^kv$. Thus eigenvectors of $T$ remain eigenvectors of every polynomial in $T$, with eigenvalues transformed by the polynomial.
\end{solution}

\begin{problem}[Zero eigenvalue and singularity]\label{prob:i09-dossier-05}
Show that $0$ is an eigenvalue of a square matrix iff the matrix is singular.
\end{problem}
\begin{solution}
$0$ is an eigenvalue iff there is $v\ne0$ with $Av=0$, i.e. iff the kernel is nontrivial. A square matrix is invertible exactly when its kernel is zero.
\end{solution}

\begin{problem}[Distinct eigenvectors]\label{prob:i09-dossier-06}
Why are eigenvectors for distinct eigenvalues independent?
\end{problem}
\begin{solution}
In a dependence relation, apply $T-\lambda_rI$ to eliminate one eigenvector while multiplying the others by nonzero scalars. Induction reduces to a shorter relation, forcing all coefficients to zero.
\end{solution}

\begin{problem}[Triangular spectrum]\label{prob:i09-dossier-07}
Read the eigenvalues of $\begin{pmatrix}3&*&*\\0&-1&*\\0&0&3\end{pmatrix}$.
\end{problem}
\begin{solution}
The characteristic polynomial of a triangular matrix is the product $(t-3)(t+1)(t-3)$, so the eigenvalues are $3$ with algebraic multiplicity $2$ and $-1$ with multiplicity $1$.
\end{solution}

\begin{problem}[Eigenspace as kernel]\label{prob:i09-dossier-08}
Why is $E_\lambda$ always a subspace?
\end{problem}
\begin{solution}
$E_\lambda=\ker(T-\lambda I)$. Kernels of linear maps are subspaces, so closure and the zero vector follow automatically.
\end{solution}

\begin{problem}[Eigenvalue of an inverse]\label{prob:i09-dossier-09}
If $T$ is invertible and $Tv=\lambda v$, what eigenvalue does $v$ have for $T^{-1}$?
\end{problem}
\begin{solution}
Invertibility implies $\lambda\ne0$. Applying $T^{-1}$ gives $v=\lambda T^{-1}v$, so $T^{-1}v=\lambda^{-1}v$.
\end{solution}

\begin{problem}[Trace and eigenvalues]\label{prob:i09-dossier-10}
For a triangularizable operator, relate trace to eigenvalues counted with algebraic multiplicity.
\end{problem}
\begin{solution}
In a triangular representation the diagonal entries are the eigenvalues with multiplicity, and trace is the sum of diagonal entries. Since trace is basis-invariant, it equals the eigenvalue sum.
\end{solution}

\begin{problem}[Determinant and eigenvalues]\label{prob:i09-dossier-11}
For a triangularizable operator, relate determinant to eigenvalues.
\end{problem}
\begin{solution}
In triangular form the determinant is the product of diagonal entries, which are the eigenvalues with algebraic multiplicity. Therefore $\det T$ equals their product.
\end{solution}

\begin{problem}[Spectral synthesis]\label{prob:i09-dossier-12}
Why can knowing all eigenvalues still be insufficient to determine an operator up to similarity?
\end{problem}
\begin{solution}
Eigenvalues record only the roots of the characteristic polynomial. Operators with the same eigenvalues can have different eigenspace dimensions or Jordan block structures, such as the identity and a nontrivial shear with eigenvalue $1$.
\end{solution}

% END VOLUME01 SOLVED DOSSIERS

\section{Exercises with complete solutions}

\begin{exercise}\label{exr:i09-01}
Find the eigenvalues of $\begin{pmatrix}2&0\\0&3\end{pmatrix}$.
\end{exercise}
\begin{hint}
First determine \(\lambda\) from \(\det(A-\lambda I)=0\), then solve \((A-\lambda I)v=0\) for the eigenspace. Form \(\det(A-\lambda I)\) first; only after finding \(\lambda\) should you solve \((A-\lambda I)v=0\).
\end{hint}
\begin{solution}
They are the diagonal entries $2$ and $3$.
\end{solution}

\begin{exercise}\label{exr:i09-02}
Find the eigenspace for eigenvalue $1$ of $\begin{pmatrix}1&1\\0&1\end{pmatrix}$.
\end{exercise}
\begin{hint}
First determine \(\lambda\) from \(\det(A-\lambda I)=0\), then solve \((A-\lambda I)v=0\) for the eigenspace. Substitute the proposed pair back into \(Av=\lambda v\) before doing any further theory.
\end{hint}
\begin{solution}
Solve $(A-I)v=0$, giving $y=0$; the eigenspace is the span of $(1,0)$.
\end{solution}

\begin{exercise}\label{exr:i09-03}
Why is the zero vector excluded from the definition of eigenvector?
\end{exercise}
\begin{hint}
First determine \(\lambda\) from \(\det(A-\lambda I)=0\), then solve \((A-\lambda I)v=0\) for the eigenspace. For a repeated eigenvalue, compute the eigenspace dimension; multiplicity in the polynomial alone does not decide diagonalizability.
\end{hint}
\begin{solution}
It satisfies $T(0)=\lambda0$ for every scalar, so including it would carry no information about the operator.
\end{solution}

\begin{exercise}\label{exr:i09-04}
Prove eigenspaces are subspaces.
\end{exercise}
\begin{hint}
First determine \(\lambda\) from \(\det(A-\lambda I)=0\), then solve \((A-\lambda I)v=0\) for the eigenspace. To prove independence, apply a polynomial in \(T\) that kills all but one eigencomponent.
\end{hint}
\begin{solution}
They are kernels of the linear maps $T-\lambda I$.
\end{solution}

\begin{exercise}\label{exr:i09-05}
Show that eigenvalues depend on the scalar field.
\end{exercise}
\begin{hint}
First determine \(\lambda\) from \(\det(A-\lambda I)=0\), then solve \((A-\lambda I)v=0\) for the eigenspace. An eigenspace is a kernel: \(E_\lambda=\ker(T-\lambda I)\). Use kernel methods to find a basis.
\end{hint}
\begin{solution}
The real rotation by $90^\circ$ has characteristic polynomial $t^2+1$, with no real roots but roots $\pm i$ over $\mathbb C$.
\end{solution}

\begin{exercise}\label{exr:i09-06}
If $T$ is invertible, can $0$ be an eigenvalue?
\end{exercise}
\begin{hint}
First determine \(\lambda\) from \(\det(A-\lambda I)=0\), then solve \((A-\lambda I)v=0\) for the eigenspace. Distinct eigenvalues give distinct invariant directions; try a linear relation and apply \(T\).
\end{hint}
\begin{solution}
No. A zero eigenvalue would give a nonzero vector in the kernel.
\end{solution}

\begin{exercise}\label{exr:i09-07}
What are the eigenvalues of a triangular matrix?
\end{exercise}
\begin{hint}
First determine \(\lambda\) from \(\det(A-\lambda I)=0\), then solve \((A-\lambda I)v=0\) for the eigenspace. Compare the sum of eigenspace dimensions with the dimension of the whole space.
\end{hint}
\begin{solution}
They are its diagonal entries, because the characteristic polynomial is the product of $t-a_{ii}$.
\end{solution}

\begin{exercise}\label{exr:i09-08}
Explain the difference between algebraic and geometric multiplicity for a shear.
\end{exercise}
\begin{hint}
First determine \(\lambda\) from \(\det(A-\lambda I)=0\), then solve \((A-\lambda I)v=0\) for the eigenspace. For triangular matrices, read candidate eigenvalues from the diagonal before computing eigenvectors.
\end{hint}
\begin{solution}
For the $2\times2$ shear, eigenvalue $1$ has algebraic multiplicity $2$ but geometric multiplicity $1$.
\end{solution}


% BEGIN VOL01-EXPANSION I09-exercises-01
\section{Graded supplementary exercises}
The following exercises extend the original exercise set without replacing it. They are grouped by mathematical role and increase in synthesis.

\subsection*{Standard computations}

\begin{exercise}[Check eigenpair]\label{exr:i09-expand-01}
Check whether \(v=(1,2)\) is an eigenvector of \(A=\begin{pmatrix}3&0\\0&3\end{pmatrix}\).
\end{exercise}
\begin{hint}
Compute \(Av\).
\end{hint}
\begin{solution}
\(Av=(3,6)=3v\), so yes, with eigenvalue \(3\).
\end{solution}

\begin{exercise}[Characteristic roots]\label{exr:i09-expand-02}
Find eigenvalues of \(\begin{pmatrix}2&1\\0&5\end{pmatrix}\).
\end{exercise}
\begin{hint}
Use triangular form.
\end{hint}
\begin{solution}
\(2\) and \(5\).
\end{solution}

\begin{exercise}[Eigenspace]\label{exr:i09-expand-03}
Find \(E_2\) for \(A=\begin{pmatrix}2&1\\0&5\end{pmatrix}\).
\end{exercise}
\begin{hint}
Solve \((A-2I)v=0\).
\end{hint}
\begin{solution}
\(y=0\), so \(E_2=\operatorname{span}(1,0)\).
\end{solution}

\begin{exercise}[Power action]\label{exr:i09-expand-04}
If \(Tv=-2v\), compute \(T^5v\).
\end{exercise}
\begin{hint}
Iterate the eigenvalue.
\end{hint}
\begin{solution}
\(T^5v=(-2)^5v=-32v\).
\end{solution}

\begin{exercise}[Inverse eigenvalue]\label{exr:i09-expand-05}
If \(T\) is invertible and \(Tv=4v\), what does \(T^{-1}\) do to \(v\)?
\end{exercise}
\begin{hint}
Apply \(T^{-1}\).
\end{hint}
\begin{solution}
\(T^{-1}v=\frac14v\).
\end{solution}

\subsection*{Proofs}

\begin{exercise}[Distinct-eigenvalue proof]\label{exr:i09-expand-06}
Prove eigenvectors for distinct eigenvalues are independent in the two-vector case.
\end{exercise}
\begin{hint}
Apply \(T-\lambda_2I\) to \(av_1+bv_2=0\).
\end{hint}
\begin{solution}
\(a(\lambda_1-\lambda_2)v_1=0\), so \(a=0\), then \(b=0\).
\end{solution}

\begin{exercise}[Eigenspace subspace]\label{exr:i09-expand-07}
Prove \(E_\lambda\) is a subspace.
\end{exercise}
\begin{hint}
Identify it as a kernel.
\end{hint}
\begin{solution}
\(E_\lambda=\ker(T-\lambda I)\), hence is a subspace.
\end{solution}

\begin{exercise}[Zero eigenvalue]\label{exr:i09-expand-08}
Prove \(0\) is an eigenvalue iff a square matrix is singular.
\end{exercise}
\begin{hint}
Translate \(Av=0\) with \(v\neq0\).
\end{hint}
\begin{solution}
A nonzero kernel vector exists exactly when the matrix is not invertible.
\end{solution}

\begin{exercise}[Polynomial action]\label{exr:i09-expand-09}
If \(Tv=\lambda v\), prove \(p(T)v=p(\lambda)v\).
\end{exercise}
\begin{hint}
Use \(T^kv=\lambda^kv\).
\end{hint}
\begin{solution}
Expand \(p\) as a sum of powers and apply linearity.
\end{solution}

\subsection*{Counterexamples}

\begin{exercise}[Counterexample: zero vector]\label{exr:i09-expand-10}
Why is the zero vector excluded from the definition of eigenvector?
\end{exercise}
\begin{hint}
Test \(T0=\lambda0\).
\end{hint}
\begin{solution}
The equation holds for every \(\lambda\), so including zero would destroy the meaning of eigenvalue.
\end{solution}

\begin{exercise}[Counterexample: repeated root]\label{exr:i09-expand-11}
Give a matrix with characteristic polynomial \((t-1)^2\) that is not diagonalizable.
\end{exercise}
\begin{hint}
Use a shear Jordan block.
\end{hint}
\begin{solution}
\(\begin{pmatrix}1&1\\0&1\end{pmatrix}\) has one-dimensional eigenspace.
\end{solution}

\begin{exercise}[Counterexample: field dependence]\label{exr:i09-expand-12}
Give a real matrix with no real eigenvalues.
\end{exercise}
\begin{hint}
Use a quarter-turn.
\end{hint}
\begin{solution}
\(\begin{pmatrix}0&-1\\1&0\end{pmatrix}\) has characteristic polynomial \(t^2+1\).
\end{solution}

\subsection*{Applications}

\begin{exercise}[Application: recurrence mode]\label{exr:i09-expand-13}
If a state starts in an eigenvector direction with eigenvalue \(0.9\), what happens after \(n\) iterations?
\end{exercise}
\begin{hint}
Use \(T^nv=\lambda^nv\).
\end{hint}
\begin{solution}
The state is scaled by \(0.9^n\) and decays to zero.
\end{solution}

\begin{exercise}[Application: unstable mode]\label{exr:i09-expand-14}
What does an eigenvalue \(1.2\) imply for repeated iteration along its eigenvector?
\end{exercise}
\begin{hint}
Compute \(1.2^n\).
\end{hint}
\begin{solution}
The magnitude grows exponentially like \(1.2^n\).
\end{solution}

\subsection*{Challenges}

\begin{exercise}[Challenge: collision parameter]\label{exr:i09-expand-15}
For \(A_t=\begin{pmatrix}t&1\\0&2\end{pmatrix}\), classify when it has two distinct eigenvalues.
\end{exercise}
\begin{hint}
Read diagonal entries.
\end{hint}
\begin{solution}
Distinct when \(t\neq2\); at \(t=2\) only eigenvalue \(2\) remains.
\end{solution}

\begin{exercise}[Challenge: trace/determinant check]\label{exr:i09-expand-16}
A \(2\times2\) matrix has eigenvalues \(3,-2\). What trace and determinant must it have?
\end{exercise}
\begin{hint}
For triangularizable matrices, use sum and product.
\end{hint}
\begin{solution}
Trace \(1\), determinant \(-6\).
\end{solution}
% END VOL01-EXPANSION I09-exercises-01
\section*{Chapter summary}
Eigenvectors are directions preserved by an operator up to scaling, and eigenvalues are the corresponding scale factors. The characteristic polynomial packages all eigenvalues algebraically. This chapter emphasizes both geometric meaning and the limitations of eigenvector methods when a full eigenbasis does not exist.
The important habit is to move fluently between invariant statements and concrete coordinates while remembering which conclusions depend on finite dimensionality, the scalar field, or the inner product.
